\section{Auswertung}
\label{sec:Auswertung}
%include tables
%\begin{table}
 \caption{Im Experiment gemessene Ströme der Sweep-Spule und der Horizontalfeldspule für die Transparenzminima beider Isotope sowie die Frequenz des angelegten RF-Wechselfelds }
 \label{tab:current}
 \centering
 \begin{tabular}{ccccc}
 \toprule 
    RF-Wechselfeld/$\si{\kilo\hertz}$ & $I_{\mathrm{Horizontal,1}}/\si{\milli \ampere}$ & $I_{\mathrm{Horizontal,2}}/\si{\milli \ampere}$ & $I_{\mathrm{Sweep,1}}/\si{\milli \ampere}$ & $I_{\mathrm{Sweep,2}}/\si{\milli \ampere}$ \\
     \midrule
     100.3 & 0 & 0 & 597 & 666 \\
     201.0 & 35 & 35 & 245 & 484 \\
     301.0 & 56 & 56 & 187 & 545 \\
     400.0 & 76 & 76 & 139 & 617 \\
     499.0 & 91 & 91 & 162 & 753 \\
     603.0 & 116 & 116 & 47 & 763 \\
     699.0 & 88 & 159 & 688 & 473 \\
     801.0 & 101 & 182 & 736 & 508 \\
     899.0 & 149 & 207 & 265 & 493 \\
     1004.0 & 156 & 237 & 410 & 424 \\
 \bottomrule
 \end{tabular}
\end{table}
\subsection{Eichen des Elektromagneten}%Tabelle mit den Messdaten für B-Feld Eichung?
Das gemessene Magnetfeld $B$ ist in Abbildung \ref{fig:eichi} gegen den Feldstrom $I$ auf dem
Intervall $\SI{1}{\ampere} \le I \le \SI{16}{\ampere}$ mit zugehörigem Ausgleichspolynom
aufgetragen. Dabei ist das Ausgleichspolynom dritter Ordnung.
Die Ausgleichsrechnung wird mit Python \cite{scipy} durchgeführt.
\begin{figure}
	\centering
	\includegraphics[width=\textwidth]{pictures/eichi.pdf}
	\caption{Gemessenes B-Feld aufgetragen gegen den Feldstrom $I$ mit zugehörigem Ausgleichspolynom.}
	\label{fig:eichi}
\end{figure}

\subsection{Vorbereitende Aufgaben}
	\label{sec:vorb}
	Bei der verwendeten Cadmium-Lampe wird die rote Linie für den normalen Zeemann-Effekt
	mit dem Übergang
	\begin{align*}
		5^1D_2 \rightarrow 5^1P_1
	\end{align*}
	erzeugt, wobei sich durch die Multiplizität $M=2S+1=1$ jeweils ein Gesamtspin von
	$S=0$ ergibt -- für den normalen Zeemann-Effekt notwendigerweise erfüllt --
	was zur Folge hat, dass der Landé-Faktor eins wird.
	Die Niveaus, die an der Erzeugung der blauen Linie beteiligt sind, tragen offensichtlich
	einen Spin, sodass der Landé-Faktor von $L$, $S$ und $J$ abhängt. Der angedeutete
	Übergang lautet
	\begin{align*}
		6^3S_1 (g_J=2) \rightarrow 5^3S_1 (g_J=\frac{3}{2}) \, \mathrm{.}
	\end{align*}

	Weiterhin wird das Dispersionsgebiet $\Delta \lambda_{\mathrm{D}}$ der
	Lummer-Gehrcke-Platte mit Formel
	\eqref{eqn:dispersionsgebiet} bestimmt. Dabei ist die Dicke der Platte $d=\SI{4}{\milli\meter}$
	und die Länge der Platte  $L = \SI{120}{\milli\meter}$. Zuletzt ergibt sich der
	Brechungsindex $n$ aus den Wellenlängen des verwendeten Lichtes zu
	$n(\lambda=\SI{644}{\nano\meter}) = 1,4567$ für das rote Licht und zu
	$n(\lambda=\SI{480}{\nano\meter}) = 1,4635$ für das blaue Licht. Damit ergibt sich
	das Dispersionsgebiet zu
	\begin{align*}
		\Delta \lambda_{\mathrm{D}}(\lambda=\SI{643.8}{\nano\meter}) &= \SI{48.913}{\pico\meter} \, \mathrm{,} \\
		\Delta \lambda_{\mathrm{D}}(\lambda=\SI{480.0}{\nano\meter}) &= \SI{26.952}{\pico\meter} \, \mathrm{.}
	\end{align*}
	Für das Auflösungsvermögen $A$ ergeben sich mit Gleichung \eqref{eqn:A} die Werte:
\begin{align*}
	A(\lambda=\SI{643.8}{\nano\meter}) &= 209063,644 \, \mathrm{,} \\
	A(\lambda=\SI{480}{\nano\meter}) &= 285458,063 \, \mathrm{.}
\end{align*}

\subsection{Bestimmung der Landé-Faktoren $g$}
	Die Landé-Faktoren werden mit Formel \eqref{eqn:g_exp} beziehungsweise \eqref{eqn:g_expnorm} bestimmt. Hierfür werden die mit der
	Digital-Kamera aufgenommenen Aufspaltungen mit und ohne B-Feld übereinander aufgetragen.
	Mit dem Bildbearbeitungsprogramm Gimp \cite{gimp} können mit der Distanzfunktion
	die benötigten
	Werte für $\Delta s$ (aus dem Bild ohne B-Feld) und $\delta s$ (aus dem Bild mit B-Feld)
	in Pixeln aus den aufgenommenen Bildern bestimmt werden. Dabei wurde jeweils die
	Strecke zwischen den intensivsten Punkten gemessen.
	Der Wert $B$ für das B-Feld wird mit dem Ausgleichspolynom bestimmt, indem die
	verwendeten Ströme eingesetzt werden.
	Es wird mit den Konstanten
	\begin{align*}
		\mathrm{h} &= 6,626 \cdot 10^{-34} \, \si{\joule\second} \, \text{, aus \cite{h}} \\
		\mathrm{c} &= \SI{299792459}{\meter\per\second}\, \text{, aus \cite{c}} \\
		\mu_{\mathrm{B}} &= 927,4 \cdot 10^{-26} \, \si{\joule\per\tesla} \, \text{, aus \cite{magneton}}
	\end{align*}
	gerechnet.
	\FloatBarrier
\subsubsection{Rote Linie}
	Zuerst ist die Zeemann-Aufspaltung mit und ohne B-Feld in Abbildung \ref{fig:roterstern}
	dargestellt.
	Die gemessenen Werte für $\delta s$ und $\Delta s$ sind in Tabelle \ref{tab:rot}
	aufgetragen. Mit dem bestimmten Dispersionsgebiet (vgl. \ref{sec:vorb}) ergibt sich
	durch Mittelwert- und Standardabweichungsbestimmung mit Python/Scipy \cite{scipy} für
	den Landé-Faktor:
	\begin{equation*}
		\alpha= g = \SI{1.014(28)} \, \text{.}
	\end{equation*}
	\begin{table}
		\centering
		\caption{Messwerte (in Pixel) zur Bestimmung des Landé-Faktors bei der roten Linie mit einem Magnetfeld von $B = \SI{652}{\milli\tesla}$.}
		\label{tab:rot}
		\begin{tabular}{c | c}
			\toprule
			$\Delta s$ & $\delta s$\\
			\midrule
			148 & 72 \\
			136 & 72 \\
			130 & 68 \\
			132 & 66 \\
			132 & 70 \\
			126 & 70 \\
			120 & 66 \\
			122 & 62 \\
			122 & 62 \\
			120 & 66 \\
			\bottomrule
		\end{tabular}
	\end{table}
\begin{figure}
	\centering
	\includegraphics[width=\textwidth]{ouurs/red_cropped.png}
	\caption{Aufgenommene Zeemann-Aufspaltung der roten Linie.}
	\label{fig:roterstern}
\end{figure}

\FloatBarrier
\subsubsection{Blaue \texorpdfstring{$\sigma-$}ALinie}
	Die Zeemann-Aufspaltung mit und ohne B-Feld ist in Abbildung \ref{fig:blau_sigma}
	zu sehen. Weiterhin sind analog wie im vorherigen Abschnitt die benötigten Messwerte
	für $\Delta s$ und $\delta s$ in Tabelle \ref{tab:blau_sigma} aufgetragen.
	Damit ergibt sich der Faktor $\alpha$, welcher die Landé-Faktoren beider Niveaus enthält, zu:
	\begin{equation*}
		\alpha = \SI{1.84(4)} \, \text{.}
	\end{equation*}
	\begin{table}
		\centering
		\caption{Messwerte (in Pixel) zur Bestimmung des Landé-Faktors bei der blauen $\sigma$-Linie mit einem Magnetfeld von $B = \SI{333}{\milli\tesla}$.}
		\label{tab:blau_sigma}
		\begin{tabular}{c | c}
			\toprule
			$\Delta s$ & $\delta s$\\
			\midrule
			126 & 60 \\
			123 & 57 \\
			117 & 61 \\
			117 & 57 \\
			114 & 60 \\
			117 & 57 \\
			111 & 54 \\
			111 & 51 \\
			112 & 57 \\
			114 & 55 \\
			\bottomrule
		\end{tabular}
	\end{table}
	\begin{figure}
	\centering
	\includegraphics[width=\textwidth]{ouurs/blau_sigma_cropped.png}
	\caption{Aufgenommene Zeemann-Aufspaltung der blauen $\sigma$-Linie.}
	\label{fig:blau_sigma}
\end{figure}

\FloatBarrier
\subsubsection{Blaue \texorpdfstring{$\pi-$}ALinie}
	Zuletzt wird der Landé-Faktor für die blaue $\pi$-Linie bestimmt. Das Vorgehen ist wieder
	analog zu den vorherigen Abschnitten.
	Die Messwerte befinden sich in Tabelle \ref{tab:blau_pi} und die Aufspaltung
	ist in Abbildung \ref{fig:blau_pi} zu finden.
	Der $\alpha$-Faktor ergibt sich hier zu:
	\begin{equation*}
		\alpha = \SI{0.488(23)} \, \text{.}
	\end{equation*}
	\begin{table}
		\centering
		\caption{Messwerte (in Pixel) zur Bestimmung des Landé-Faktors bei der blauen $\pi$-Linie mit einem Magnetfeld von $B = \SI{1001}{\milli\tesla}$.}
		\label{tab:blau_pi}
		\begin{tabular}{c | c}
			\toprule
			$\Delta s$ & $\delta s$\\
			\midrule
			126 & 48 \\
			120 & 51 \\
			111 & 45 \\
			117 & 46 \\
			111 & 45 \\
			105 & 48 \\
			108 & 36 \\
			102 & 42 \\
			108 & 36 \\
			100 & 35 \\
			\bottomrule
		\end{tabular}
	\end{table}
	\begin{figure}
	\centering
	\includegraphics[width=\textwidth]{ouurs/blau_pi_crop.png}
	\caption{Aufgenommene Zeemann-Aufspaltung der blauen $\pi$-Linie.}
	\label{fig:blau_pi}
\end{figure}
